% Altimeter Pipeline — Bed Level Measurement Methods Section
% Drop into your paper's methods section. Adjust \subsection depth as needed.
% Requires: amsmath, siunitx (optional, for units), natbib or biblatex

\subsection{Bed level measurements}

\subsubsection{Instrumentation}

Bed level change was measured using bottom-mounted acoustic instruments
manufactured by Echologger (EofE Ultrasonics Co., Ltd), deployed on
vertical pipes with the transducer facing downward toward the seabed.
Two instrument types were used:

\begin{itemize}
  \item \textbf{AA400 altimeter}: A \SI{450}{\kilo\hertz} autonomous
    acoustic altimeter that measures the distance from the sensor to the
    nearest acoustically reflective surface (the seabed). The instrument
    records altitude, temperature, battery voltage, and acoustic return
    amplitude at \SI{2}{\hertz}. Data are stored as ASCII CSV files on
    internal flash memory, with deployment durations limited to
    approximately 30~days by memory capacity.

  \item \textbf{EA400 echosounder}: A \SI{450}{\kilo\hertz} autonomous
    precision echosounder that measures both the distance to the seabed
    and a full acoustic backscatter intensity profile through the water
    column. The backscatter profile is resolved into depth bins of
    \SI{7.5}{\milli\meter} spacing over a \SIrange{2}{3}{\meter} range,
    with the number of bins varying by instrument configuration (272 or
    400~bins). The instrument additionally records pitch, roll, and
    temperature. Data are stored in a proprietary binary format (.BIN)
    on SD card, with deployment durations limited to approximately
    30~days by battery capacity. The binary file format was
    reverse-engineered for this study (Appendix~\ref{app:bin_format}).
\end{itemize}

\subsubsection{Deployment configuration}

Instruments were deployed at three sites along the San Diego County
coastline: south of the SIO pier on MOP~511 at \SI{6}{\meter} depth,
along MOP~586 at Torrey Pines State Beach at \SI{5}{\meter},
\SI{7}{\meter}, \SI{10}{\meter}, and \SI{15}{\meter} depths, and at
Solana Beach on MOP~654 at \SI{7}{\meter} depth.

At the SIO~pier site, both an AA400 altimeter and EA400 echosounder
were co-deployed and swapped monthly. At Torrey Pines, AA400 altimeters
were deployed from February--August 2024 and were subsequently replaced
by EA400 echosounders at the \SI{5}{\meter}, \SI{7}{\meter}, and
\SI{10}{\meter} locations, while the \SI{10}{\meter} and
\SI{15}{\meter} locations retained AA400 altimeters from November 2024
onward. At Solana Beach, only EA400 echosounders were deployed.

The AA400 altimeters sampled in burst mode, recording 60--1800 samples
(30~seconds to 15~minutes) at \SI{2}{\hertz} every 20--30~minutes,
depending on site and deployment period. The EA400 echosounders
sampled similarly in burst mode, typically recording 300~samples
(2.5~minutes) at \SI{2}{\hertz} every 30~minutes.

Most altimeter and echosounder deployments were co-located with Nortek
Vector ADVs configured in pressure--velocity (PUV) mode, providing
simultaneous measurements of wave height, period, direction, orbital
velocity, and bed shear stress.

\subsection{Level 1: Data reading and concatenation}

\subsubsection{AA400 altimeter}

Raw AA400 data were read from Nortek RangeLogger ASCII log files. Each
file contains a device configuration header (instrument settings,
serial number, sampling parameters) followed by comma-delimited data
rows with timestamp, ping count, altitude, temperature, battery
voltage, and acoustic return amplitude. A burst summary line appended
to the last sample of each burst (containing sample count and interval
metadata) was detected and excluded during parsing. Non-data lines
(device shutdown messages) at the end of the file were filtered by
rejecting rows whose first character was non-numeric.

\subsubsection{EA400 echosounder}

Raw EA400 binary (.BIN) files were read using a custom parser
developed for this study. Each file consists of a 128-byte header
containing instrument configuration (sound speed, timer settings)
followed by repeating records, one per ping. Each record contains a
64-byte DATA sub-record (ping number, Unix timestamp, and the
backscatter intensity profile as $N$ unsigned 16-bit integers, where
$N$ is the number of depth bins), followed by a 64-byte STAT
sub-record containing the authoritative per-ping metadata:
temperature (\si{\celsius}), altitude (\si{\meter}), pitch
(\si{\degree}), and roll (\si{\degree}). The full file was read into
memory as a byte vector and reshaped into a record matrix for
vectorized extraction of all fields, enabling efficient processing
of files exceeding \SI{500}{\mega\byte}.

Multiple files spanning a single deployment were concatenated and
sorted chronologically.

\subsection{Level 2: Quality control}

\subsubsection{Despiking}

Acoustic altimeter data contain impulsive spikes caused by
intermittent acoustic interference, biological targets (fish, kelp),
and multipath reflections. Spikes were identified and removed using
the phase-space threshold method of \citet{goring2002}, as modified
by \citet{mori2005}. This method constructs an ellipsoid in the
three-dimensional phase space defined by the signal $x$, its first
derivative $dx/dt$, and second derivative $d^2x/dt^2$. The ellipsoid
semi-axes are set to $\lambda \sigma$ for each dimension, where
$\sigma$ is the standard deviation and $\lambda = \sqrt{2 \ln N}$ is
the Universal threshold. Points lying outside the ellipsoid are
classified as spikes and replaced with NaN. The algorithm iterates
(up to 20 iterations) until no additional spikes are detected. This
approach is parameter-free (the threshold derives from the data
itself) and has been widely validated for acoustic Doppler velocimeter
data \citep{goring2002}.

Prior to phase-space despiking, zero-altitude readings (indicating no
acoustic return) were flagged and set to NaN. Quality flags were
stored as a uint16 bitmask: bit~1 for invalid/zero values, bit~2 for
phase-space spikes.

\subsubsection{Echosounder tilt correction and masking}

For EA400 echosounder deployments, a geometric tilt correction was
applied to account for the overestimation of distance-to-bed caused
by non-vertical beam orientation. A sensor tilted by angle $\alpha$
from vertical measures an acoustic path length $r = d / \cos\alpha$,
where $d$ is the true vertical distance. The per-ping total tilt angle
was computed as $\alpha = \sqrt{\theta_\text{pitch}^2 +
\theta_\text{roll}^2}$, and the corrected altitude was obtained as:
%
\begin{equation}
  a_\text{corrected} = a_\text{measured} \cos\alpha
\end{equation}
%
For the Solana Beach deployment (baseline tilt
$\alpha = \SI{13.2}{\degree}$), the correction reduced the measured
altitude by \SI{24}{\milli\meter} at the median range of
\SI{910}{\milli\meter} (a 2.6\% correction). For instruments with
baseline tilt $< \SI{2}{\degree}$ (Torrey Pines, SIO pier), the
correction is $< \SI{1}{\milli\meter}$ and effectively negligible.
Because the tilt is stable over time, the correction is a constant
offset and does not affect relative bed level change $\Delta z$; it
is relevant only for absolute elevation comparisons.

The pitch and roll sensors were also used to identify periods when
the instrument was disturbed (e.g., by wave loading on kelp-fouled
pipes). Rather than masking based on absolute tilt---which would
reject all data from instruments installed with a static offset from
vertical---the tilt mask flags deviations from the baseline
(deployment median) pitch and roll. Pings where
$|\theta_\text{pitch} - \tilde{\theta}_\text{pitch}| > \SI{2}{\degree}$
or
$|\theta_\text{roll} - \tilde{\theta}_\text{roll}| > \SI{2}{\degree}$
were flagged, and both altitude and backscatter data at those
timestamps were set to NaN. This deviation-based approach preserves
data from instruments with static tilt offsets up to
\SI{13}{\degree} (observed at the Solana Beach site).

\subsubsection{Backscatter masking}

The echosounder backscatter profile extends from the sensor face to
its maximum configured range (\SIrange{2}{3}{\meter}). Acoustic
returns from depth bins beyond the detected seabed represent
sub-surface reflections and are not physically meaningful for water
column analysis. For each ping, all backscatter bins at ranges
exceeding the detected altitude were set to NaN. Pings with no
detected bed echo (altitude = 0 or NaN) had their entire backscatter
profile set to NaN, as the absence of a bed return indicates
unreliable acoustic conditions for the full profile.

\subsection{Level 3: Derived products}

\subsubsection{Burst-averaged bed level}

The 2~Hz altitude time series was reduced to burst-averaged bed level
change matched to the temporal resolution of the co-located PUV wave
measurements. The burst structure was auto-detected from the time
series: instruments with time gaps exceeding \SI{10}{\second} between
consecutive samples (indicating burst-mode acquisition) were processed
using the physical burst boundaries, while continuously sampled
instruments were divided into non-overlapping 2048-sample
(\SI{17.1}{\minute}) windows matching the PUV spectral segment length.

For each burst window, the \emph{median} altitude was computed,
providing a measure robust to both wave-frequency apparent bed
oscillations (caused by orbital motion modulating the acoustic path
length) and any remaining undetected spikes. The interquartile range
(IQR) of altitude within each burst was recorded as an estimate of
measurement uncertainty, encompassing both wave-orbital contamination
and acoustic noise. Bursts with fewer than 50\% valid (non-NaN)
samples were rejected.

Bed level change was computed as:
%
\begin{equation}
  \Delta z(t_k) = -(\ \tilde{a}_k - \tilde{a}_0\ )
\end{equation}
%
where $\tilde{a}_k$ is the burst-median altitude at time $t_k$ and
$\tilde{a}_0$ is the median altitude of the first fully valid burst.
The sign convention is: accretion (bed moving upward, reducing the
sensor-to-bed distance) is positive; erosion is negative.

\subsubsection{Bed level change rate}

The bed level change rate between consecutive bursts was computed as:
%
\begin{equation}
  \frac{d(\Delta z)}{dt}\bigg|_{k} =
  \frac{\Delta z(t_k) - \Delta z(t_{k-1})}{t_k - t_{k-1}}
\end{equation}
%
expressed in \si{\milli\meter\per\hour}. Rates computed across time
gaps exceeding three times the median burst interval were set to NaN
to avoid spurious values across data outages.

\subsubsection{Burst-averaged backscatter}

For EA400 echosounder deployments, the backscatter profile was
averaged over each burst window to produce burst-mean and burst-maximum
profiles. The burst-maximum profile captures peak suspension events
within each averaging window and is useful for identifying
wave-group-driven sediment mobilization.

\subsection{Validation}

The pipeline was validated across five representative deployments
spanning all three sites and both instrument types
(Table~\ref{tab:alt_validation}):

\begin{enumerate}
  \item SIO pier AA400 altimeter (continuous mode, 6m): 257~bursts of
    $\sim$1800 samples, median IQR = \SI{4.6}{\milli\meter}.
  \item SIO pier AA400 altimeter (burst mode, 6m): 1531~bursts of
    $\sim$300 samples, median IQR = \SI{5.1}{\milli\meter}.
  \item Torrey Pines EA400 echosounder (10m): 1719~bursts, median
    IQR = \SI{7.2}{\milli\meter}, 4\% rejected.
  \item Torrey Pines AA400 altimeter (burst mode, 10m): 6883~bursts
    of $\sim$60 samples, median IQR = \SI{3.8}{\milli\meter}.
  \item Solana Beach EA400 echosounder (7m, tilted pipe): 2600~bursts,
    19\% rejected by tilt deviation mask, median
    IQR = \SI{7.3}{\milli\meter}.
\end{enumerate}

The AA400 altimeters consistently exhibited lower within-burst IQR
(\SIrange{3.8}{5.1}{\milli\meter}) compared to the EA400 echosounders
(\SIrange{7.2}{7.3}{\milli\meter}), reflecting the simpler measurement
geometry and more focused acoustic beam of the dedicated altimeter.
The percent of valid samples within accepted bursts exceeded 88\%
across all deployments. The Solana Beach deployment demonstrated
that the deviation-based tilt mask correctly preserves data from a
tilted instrument (baseline tilt \SI{13.2}{\degree}) while rejecting
periods of transient disturbance.

\subsubsection{Independent validation against beach surveys}

Echosounder bed level measurements were independently validated
against beach morphology surveys conducted by the Scripps Coastal
Processes Group. Subaqueous profiles were obtained at approximately
monthly intervals using a jetski-mounted RTK GPS system, reaching
water depths of \SIrange{8}{11}{\meter} NAVD88 at MOP~654 (Solana
Beach). Survey profiles were interpolated to the instrument cross-shore
position (identified by converting the instrument latitude/longitude
to the MOP coordinate system) to extract the surveyed bed elevation
at each survey date. The bed level \emph{change} between consecutive
survey pairs ($\Delta z_\text{survey}$) was compared against the
corresponding altimeter-measured change ($\Delta z_\text{echo}$)
over the same interval.

For the Solana Beach EA400 echosounder at \SI{7}{\meter} depth over
the period March--June 2025, four jetski GPS surveys reached the
instrument depth, yielding two paired $\Delta z$ comparisons. The
RMSE of the residual ($\Delta z_\text{echo} - \Delta z_\text{survey}$)
was \SI{22}{\milli\meter} with a bias of \SI{-4}{\milli\meter}. The
echosounder correctly captured the sign and approximate magnitude
of bed level changes observed in the surveys.

The \SI{22}{\milli\meter} RMSE is comparable to the spatial variability
inherent in the survey comparison itself: survey-to-survey elevation
changes at MOP~654 had a standard deviation of \SI{21}{\milli\meter},
and the neighboring MOP~655 transect (\SI{100}{\meter} alongshore)
showed \SI{37}{\milli\meter} standard deviation at the same
cross-shore position. This spatial aliasing arises because the jetski
survey line does not pass directly over the instrument, and bedforms
(ripples, megaripples) cause centimeter-scale elevation variability
over meter-scale horizontal distances. The validation RMSE is therefore
at the noise floor of the survey comparison, indicating that the
echosounder accuracy is at least as good as can be resolved by the
available ground truth.

% TODO: Insert Table \ref{tab:alt_validation} with deployment summary
% statistics (N_bursts, N_rejected, BL_range, IQR_med, pctValid)
